%!TEX program = xelatex
\def\ano{2023}
\def\semestre{1}
\def\disciplina{Introdução à Álgebra Linear}
\def\turma{3}
\def\autor{Jos\'e Ant\^onio O. Freitas}
\def\instituto{MAT-UnB}

\documentclass{beamer}
\usetheme{Madrid}
\usecolortheme{beaver}
% \mode<presentation>
\usepackage[singlelinecheck=false]{caption}
\usepackage{textpos}
\usepackage{amssymb}
\usepackage{amsmath,amsfonts,amsthm,amstext}
\usepackage[justification=centering]{caption}
\usepackage[brazil]{babel}
% \usepackage[latin1]{inputenc}
\usepackage{graphicx}
\graphicspath{{../../figuras_geogebra/}{../../Pictures/}}
\usepackage{enumitem}
\usepackage{multicol}
\usepackage{answers}
\usepackage{qrcode}

\usepackage{pifont}
\newcommand{\cmark}{\ding{228}}%

% requires version 0.3 of the package
\usepackage[customcolors,beamer]{hf-tikz}

\tikzset{style green/.style={
        set fill color=green!50!lime!60,
        set border color=white,
    },
    style cyan/.style={
        set fill color=cyan!90!blue!60,
        set border color=white,
    },
    style orange/.style={
        set fill color=orange!80!red!60,
        set border color=white,
    },
    hor/.style={
        above left offset={-0.15,0.31},
        below right offset={0.15,-0.125},
        #1
    },
    ver/.style={
        above left offset={-0.1,0.3},
        below right offset={0.15,-0.15},
        #1
    }
}

\usepackage{tikz,ifthen,tikz-3dplot}
\usepackage{siunitx}
\usetikzlibrary{circuits.ee.IEC}
\usetikzlibrary{lindenmayersystems}
\usetikzlibrary[shadings]
\usepackage{pgfplots}
\pgfplotsset{compat=1.15}
\usepackage{mathrsfs}
\usetikzlibrary{arrows}
\usepackage{tcolorbox}
\newcounter{saveenumi}
\newcommand{\seti}{\setcounter{saveenumi}{\value{enumi}}}
\newcommand{\conti}{\setcounter{enumi}{\value{saveenumi}}}

\resetcounteronoverlays{saveenumi}


\renewcommand{\qedsymbol}{$\blacksquare$}

%\setbeamertemplate{itemize items}[ball]

%\setbeamertemplate{itemize items}[default]
%-------------amatrix
% Augmented matrix.  Usage (note the argument does not count the aug col):
% \begin{amatrix}{2}
    %   1  2  3 \\  4  5  6
    % \end{amatrix}
\newenvironment{amatrix}[1]{%
    \left[\begin{array}{@{}*{#1}{r}|r@{}}
    }{%
    \end{array}\right]
}

\newtheorem{definicao}{Defini\c{c}\~ao}[section]
\newtheorem{definicoes}{Defini\c{c}\~oes}[section]
\newtheorem{exemplo}{Exemplo}[section]
\newtheorem{exemplos}{Exemplos}[section]
\newtheorem{exercicio}{Exerc{\'\i}cio}
\newtheorem{observacao}{Observa{\c c}{\~a}o:}[section]
\newtheorem{observacoes}{Observa{\c c}{\~o}es:}[section]
\newtheorem*{solucao}{Solu{\c c}{\~a}o:}
\newtheorem{proposicao}{Proposi\c{c}\~ao}
\newtheorem{lema}{Lema}
\newtheorem{teorema}{Teorema}
\newtheorem{corolario}{Corol\'ario}
\newenvironment{prova}[1][Prova]{\noindent\textbf{#1:} }{\qedsymbol}%{\ \rule{0.5em}{0.5em}}
\def\ano{2023}
\def\semestre{2}
\def\turma{11}
\newcommand{\nsub}{\varsubsetneq}
\newcommand{\vaz}{\emptyset}
\newcommand{\dom}{{\rm dom\,}}
\newcommand{\im}{{\rm Im\,}}
\newcommand{\sub}{\subseteq}
\newcommand{\n}{\mathbb{N}}
\newcommand{\z}{\mathbb{Z}}
\newcommand{\rac}{\mathbb{Q}}
\newcommand{\real}{\mathbb{R}}
\newcommand{\complex}{\mathbb{C}}
\newcommand{\cp}[1]{\mathbb{#1}}
\newcommand{\ch}{\mbox{\textrm{car\,}}\nobreak}
\newcommand{\p}[1]{\mbox{\textrm{posto(#1)}}}
\newcommand{\vesp}[1]{\vspace{ #1  cm}}
\newcommand{\compcent}[1]{\vcenter{\hbox{$#1\circ$}}}
\newcommand{\comp}{\mathbin{\mathchoice
{\compcent\scriptstyle}{\compcent\scriptstyle}
{\compcent\scriptscriptstyle}{\compcent\scriptscriptstyle}}}
\newcommand\T{\rule{0pt}{2.6ex}}

\title{Sistemas Lineares}
\author[\autor]{\autor}
\institute[\instituto]{\instituto}
\date{}

\begin{document}
    \begin{frame}
        \maketitle
    \end{frame}

    \logo{\includegraphics[scale=.1]{logo-MAT.png}\vspace*{8.5cm}}

    \begin{frame}
        Seja $\cp{K} = \rac$, $\real$ ou $\complex$.\pause

        \vspace{.3cm}

        Vamos trabalhar com \textbf{equações lineares} \pause em $q$ variáveis $x_1$, $x_2$, \dots, $x_q$ \pause que são equações que podem ser escritas na forma
        \begin{equation}\label{equacao_linear}
            a_1x_1 \pause + a_2x_2 + \pause \cdots + a_qx_q \pause = b
        \end{equation}

        onde $a_1$, $a_2$, \dots, $a_q$, \pause $b$ são escalares no corpo $\cp{K}$, \pause sendo que nem todos os coeficientes $a_i$ são nulos.
        
        \vspace{.3cm}

        No caso em que $b = 0$, \pause a equação \eqref{equacao_linear} tem a forma\pause
        \begin{equation}
            a_1x_1 + a_2x_2 + \cdots + a_qx_q = 0\pause
        \end{equation}
        e é chamada de uma \textbf{equação linear homogênea}.
        
    \end{frame}

    \begin{frame}
        Seja $\cp{K} = \real$.\pause

        \vspace{.3cm}
        
        Dado um conjunto de equações da forma
        \begin{align}
            \begin{cases}\label{sistema_linear_2x2}
                ax + by = b_1\\
                cx + dy = b_2
            \end{cases}
        \end{align}
        onde $a$, $b$, $c$, $d$, $b_1$, $b_2 \in \cp{K}$.\pause
    
        \vspace{.3cm}

        Existem valores de $x$, $y \in \cp{K}$ que satisfazem as duas equações simultaneamente?
    \end{frame}

    \begin{frame}
        \definecolor{qqqqff}{rgb}{0,0,1}
        \definecolor{ffqqqq}{rgb}{1,0,0}
        \begin{center}
            \begin{tikzpicture}[line cap=round,line join=round,>=triangle 45,x=1cm,y=1cm]
                \begin{axis}[
                    x=1cm,y=1cm,
                    axis lines=middle,
                    xticklabels={,,},
                    yticklabels={,,},
                    xlabel=$x$,
                    ylabel=$y$,
                    xmin=-2,
                    xmax=5,
                    ymin=-1.3,
                    ymax=3,]
                    \draw [line width=5.2pt,dash pattern=on 1pt off 1pt,color=ffqqqq,domain=-1.348257839721257:5.807994044102577] plot(\x,{(--4-1*\x)/2});
                    \draw [line width=2pt,color=qqqqff,domain=-1.348257839721257:5.807994044102577] plot(\x,{(--8-2*\x)/4});
                    \begin{scriptsize}
                        \draw[color=ffqqqq] (1.8,2) node {$a_{11}x + a_{12}y = b_1$};
                        \draw[color=qqqqff] (3,1.5) node {$a_{21}x + a_{22}y = b_2$};
                    \end{scriptsize}
                \end{axis}
            \end{tikzpicture}
        \end{center}

        Existem infinitos valores de $x$, $y \in \cp{K}$ \pause que satisfazem as duas equações simultanemente. \pause Neste caso dizemos que o sistema \eqref{sistema_linear_2x2} \pause é \textbf{possível e indeterminado}.\pause
    \end{frame}

    \begin{frame}
        \definecolor{ffqqqq}{rgb}{1,0,0}
        \definecolor{qqqqff}{rgb}{0,0,1}
        \begin{center}
            \begin{tikzpicture}[line cap=round,line join=round,>=triangle 45,x=1cm,y=1cm]
                \begin{axis}[
                    x=1cm,y=1cm,
                    axis lines=middle,
                    xticklabels={,,},
                    yticklabels={,,},
                    xlabel=$x$,
                    ylabel=$y$,
                    xmin=-5.3,
                    xmax=6.3,
                    ymin=-1,
                    ymax=4.3,]
                    \draw [line width=2pt,color=qqqqff,domain=-5:6] plot(\x,{(--4-1*\x)/5});
                    \draw [line width=2pt,color=ffqqqq,domain=-5:6] plot(\x,{(--16-1*\x)/5});
                    \begin{scriptsize}
                        \draw[color=qqqqff] (-3.8,2.2) node {$a_{21}x + a_{22}y = b_2$};
                        \draw[color=ffqqqq] (-3.8,3.4) node {$a_{11}x + a_{12}y = b_1$};
                    \end{scriptsize}
                \end{axis}
            \end{tikzpicture}
        \end{center}

        Não existem valores de $x$, $y \in \cp{K}$ \pause satisfazendo as duas equações simultaneamente. \pause Nesse caso dizemos que o sistema \eqref{sistema_linear_2x2} \pause é \textbf{impossível}.

    \end{frame}

    \begin{frame}
        \definecolor{qqqqff}{rgb}{0,0,1}
        \definecolor{ffqqqq}{rgb}{1,0,0}
        \begin{center}
            \begin{tikzpicture}[line cap=round,line join=round,>=triangle 45,x=1cm,y=1cm]
                \begin{axis}[
                    x=1cm,y=1cm,
                    axis lines=middle,
                    xticklabels={,,},
                    yticklabels={,,},
                    xlabel=$x$,
                    ylabel=$y$,
                    xmin=-2,
                    xmax=4.3,
                    ymin=-2,
                    ymax=3.5,]
                    \draw [line width=2pt,color=ffqqqq,domain=-4:5] plot(\x,{(--4.94--7.92*\x)/7.58});
                    \draw [line width=2pt,color=qqqqff,domain=-4:5] plot(\x,{(--4.824-8.66*\x)/8.24});
                    \begin{scriptsize}
                        \draw[color=ffqqqq] (2.9,2) node {$a_{11}x + a_{12}y = b_1$};
                        \draw[color=qqqqff] (2.7,-0.5) node {$a_{21}x + a_{22}y = b_2$};
                    \end{scriptsize}
                \end{axis}
            \end{tikzpicture}
        \end{center}
        
        Nesca caso existe um único conjunto de valores $x = \alpha \in \cp{K}$ e $y = \lambda \in \cp{K}$ \pause que satisfaz as duas equações simultaneamente. \pause Neste caso dizemos que o sistema \eqref{sistema_linear_2x2} \pause é \textbf{possível e determinado}.\pause
    \end{frame}

    \begin{frame}
        Esse mesmo tipo de análise pode ser aplicada num sistema da forma
        \begin{align}
            \begin{cases}\label{sistema_linear_3x3}
                a_1x + a_2y + a_3z = b_1\\
                a_4x + a_5y + a_6z = b_2\\
                a_7x + a_8y + a_9z = b_3
            \end{cases}
        \end{align}
        onde $a_i \in \cp{K}$ para $1 \le i \le 9$ e $b_j \in \cp{K}$ para $1 \le j \le 3$.\pause

        \vspace{.3cm}

        Existem valores de $x$, $y$ e $z \in \cp{K}$ que satisfaçam as três equações simultaneamente?
    \end{frame}
    
    \begin{frame}
        Como no caso anterior, temos as seguintes possibilidades:\pause
        \begin{enumerate}[label={\roman*})]
            \item Existe um único conjunto de valores $x = \alpha \in \cp{K}$, $y = \lambda \in \cp{K}$ e $z = \gamma \in \cp{K}$ \pause que satisfaz as três equações simultaneamente. \pause Neste caso dizemos que o sistema \eqref{sistema_linear_3x3} \pause é \textbf{possível e determinado}.\pause

            \item Existem infinitos valores de $x$, $y$, $z \in \cp{K}$ \pause que satisfazem as três equações simultanemente. \pause Neste caso dizemos que o sistema \eqref{sistema_linear_3x3} \pause é \textbf{possível e indeterminado}.\pause

            \item Não existem valores de $x$, $y$, $z \in \cp{K}$ \pause satisfazendo as três equações simultaneamente. \pause Nesse caso dizemos que o sistema \eqref{sistema_linear_3x3} \pause é \textbf{impossível}.
        \end{enumerate}
    \end{frame}

    \begin{frame}

        Escolha escalares $b_1$, \dots, $b_m$ \pause e $a_{ij}$, \pause $1 \le i \le m$, $1 \le j \le n$ todos em $\cp{K}$.

        \vspace{.3cm}

        Queremos saber se é possível encontrar valores para  $x_1$, $x_2$, \dots, $x_n$ \pause de modo que o seguinte conjunto de equações sejam válidas: \pause
        \begin{equation}\label{sistema_linear_geral}
	    \begin{cases}
                a_{11}x_1 + a_{12}x_2 + \cdots + a_{1n}x_n = b_1\\\pause
                a_{21}x_1 + a_{22}x_2 + \cdots + a_{2n}x_n = b_2\\\pause
                \qquad \vdots\\
                a_{m1}x_1 + a_{m2}x_2 + \cdots + a_{mn}x_n = b_m\pause
            \end{cases}
        \end{equation}
        
        O conjunto de equações em \eqref{sistema_linear_geral} é chamado de um \pause \textbf{sistema de $m$ equa\c{c}\~oes lineares \pause a $n$ inc\'ognitas} $x_1$, $x_2$, \dots, $x_n$\pause , ou simplesmente de um \textbf{sistema linear}, nas incógnitas \pause $x_1$, $x_2$, \dots, $x_n$.\pause
    \end{frame}

    \begin{frame}
        \vspace{.3cm}

        Uma solução de um sistema linear do tipo \eqref{sistema_linear_geral}\pause
        \[
            x_1 = \alpha_1, \pause x_2 = \alpha_2, \pause \dots, x_n = \alpha_n\pause
        \]
        onde $\alpha_1$, $\alpha_2$, \dots, $\alpha_n \in \cp{K}$, \pause pode ser escrita como\pause
        \[
            (\alpha_1, \alpha_2, \dots, \alpha_n)\pause
        \]
        e é chamada de uma \textbf{ênupla ordenada} ou uma \textbf{n-upla ordenada}.
    \end{frame}

    \begin{frame}
        Se $b_1 = b_2 = \cdots = b_m = 0_\cp{K} \in K$, \pause dizemos que o sistema\pause
        \begin{equation}\label{sistemalinearhomogeneo}\index{Sistema Linear}
            \begin{cases}
                a_{11}x_1 + a_{12}x_2 + \cdots + a_{1n}x_n = 0_\cp{K}\\\pause
                a_{21}x_1 + a_{22}x_2 + \cdots + a_{2n}x_n = 0_\cp{K}\\\pause
                \qquad \vdots\\
                a_{m1}x_1 + a_{m2}x_2 + \cdots + a_{mn}x_n = 0_\cp{K}\pause
            \end{cases}
        \end{equation}
        \'e um \textbf{sistema linear homog\^eneo}. \pause

        \vspace{.3cm}

        Observe que tal sistema sempre possui solu\c{c}\~ao, \pause a saber, $x_1 = x_2 = \cdots = x_q = 0_\cp{K}$.
    \end{frame}

    \begin{frame}
        No caso de um sistema linear da forma \eqref{sistema_linear_geral}, \pause o processo para encontrar suas soluções ser\'a feito mediante o uso de 3 tipos de opera\c{c}\~oes. \pause S\~ao elas:\pause
        \begin{itemize}
	    \item[$e_1$)] Troca da posi\c{c}\~ao de duas equa\c{c}\~oes.\pause
	    \item[$e_2$)] Multiplica\c{c}\~ao de uma equa\c{c}\~ao por um escalar n\~ao nulo.\pause
	    \item[$e_3$)] Substitui\c{c}\~ao de uma equa\c{c}\~ao pela soma desta equa\c{c}\~ao com alguma outra.\pause
        \end{itemize}

        Estas tr\^es opera\c{c}\~oes s\~ao chamadas de \pause \textbf{opera\c{c}\~oes elementares}.
    \end{frame}
\end{document}
